% Conclusion section drafted from existing paper sections.
% This file is a section fragment, not a standalone LaTeX document.

\section{Conclusion}
\label{sec:conclusion}

This paper studied low-bit post-training quantization for a CIFAR-10 convolutional neural network, with a focus on INT4 activation quantization. Starting from a trained FP32 CIFAR-style ResNet-18 model, we evaluated simulated PTQ configurations including INT8-MinMax, INT4-MinMax, INT4-P99.9, and the proposed INT4-MSE-Selected method. The proposed method performs calibration-time layer-wise activation clipping: for each observed post-ReLU activation site, it searches over a predefined percentile candidate set and selects the threshold that minimizes activation reconstruction MSE on the calibration set.

The experimental results show that plain INT4-MinMax suffers from a clear accuracy degradation compared with FP32 and INT8-MinMax. Introducing activation clipping substantially reduces this degradation. In particular, INT4-MSE-Selected improves Top-1 accuracy over INT4-MinMax and reduces both activation reconstruction MSE and logit MSE. The layer-wise visualization suggests that different activation sites may prefer different clipping percentiles within the predefined candidate set, indicating that activation quantization behavior can vary across layers in this controlled run. At the same time, INT4-MSE-Selected does not outperform the fixed INT4-P99.9 baseline in Top-1 accuracy in this experiment. Therefore, the results should be interpreted conservatively: the proposed method alleviates the degradation of INT4-MinMax and reduces reconstruction error, but it should not be claimed as a universally superior clipping strategy.

This work has several limitations. First, the experiments are conducted only on CIFAR-10. Second, the evaluation uses only one compact CNN setting, namely a CIFAR-style ResNet-18, so the conclusions may not directly generalize to other network architectures. Third, the reported results come from one random seed and one experimental run, so the observed differences should not be interpreted as statistically validated performance differences. Fourth, the implementation is simulated PTQ based on fake quantization and dequantization, rather than a real integer inference implementation. Fifth, no real INT4 hardware latency, throughput, energy, or deployment speed measurement is reported. Sixth, no quantization-aware training is performed, so this work only evaluates the post-training quantization setting. Finally, no comparison with stronger reconstruction-based PTQ methods such as AdaRound or BRECQ is included, so the results should be viewed as a controlled comparison among the implemented baselines rather than a comprehensive PTQ benchmark.

Future work will extend the evaluation to MobileNetV2 and other lightweight CNN architectures to test whether the observed layer-wise clipping behavior is consistent across models. Another direction is mixed-precision quantization, where different layers may use different bit-widths according to sensitivity or hardware constraints. Hardware-aware quantization is also important for connecting simulated quantization metrics with actual deployment cost. Finally, deploying the quantized CNN pipeline on an RFSoC platform would allow future experiments to measure real inference latency, resource usage, and hardware behavior under practical edge-deployment constraints.
